% related_work.tex
\section{Related Work}
\label{sec:related}

\subsection{Architectural Progress on Pedestrian Crossing Intention}
\label{sec:related-arch}

PIE~\cite{rasouli2019pie} and JAAD~\cite{rasouli2017jaad} have been the
dominant benchmarks for pedestrian crossing-intention prediction since
their introduction, and a substantial body of work has advanced
in-dataset accuracy on them through progressively richer cross-modal
fusion: SF-GRU~\cite{rasouli2020sfgru} stacks per-modality GRUs;
PCPA~\cite{kotseruba2021pcpa} adds temporal and modality attention over
3D-convolutional, pose, and box streams; more recent work replaces
recurrent fusion with transformer- or state-space-based temporal
encoders and heavier visual backbones -- e.g.\
Rahman~et~al.~\cite{adapt2026} combine a weight-shared Swin
Transformer~V2 backbone with a Mamba-based motion-encoding module and
sparse cross-modal attention (17.23\,ms/sample inference; 0.90 AUC on
PIE), and TrajFusionNet~\cite{trajfusionnet2026} fuses a sequential
branch (trajectory and ego-vehicle speed) with a visual branch
(predicted pedestrian bounding boxes overlaid on scene imagery) through
a transformer-based architecture. This line of work reports steadily
improving in-dataset accuracy and AUC on PIE and JAAD, but -- to the
best of our knowledge -- none of it evaluates cross-dataset
generalization; every result we found in this literature is reported
under the standard train-test-same-dataset protocol, confirmed by
direct inspection of released code where available (e.g.\ the ADAPT
implementation trains and evaluates PIE and JAAD independently, with no
cross-dataset transfer).

\subsection{Efficiency-Focused Approaches}
\label{sec:related-efficiency}

A parallel line of work optimizes for inference cost rather than
peak accuracy. EfficientPIE~\cite{efficientpie2025} argues that a
single observed frame, processed by a lightweight convolutional
architecture (Fused-MBConv / MBConv blocks with depthwise-separable
convolutions), is sufficient for competitive crossing-intention
accuracy, reporting 0.21\,ms inference latency and state-of-the-art
in-dataset accuracy on both PIE (0.92) and JAAD (0.89) -- outperforming
SF-GRU and every other baseline in their comparison table. This
demonstrates that architectural efficiency and competitive in-dataset
accuracy are not in tension for this task. However, EfficientPIE's
evaluation, like the architectures in
Section~\ref{sec:related-arch}, is entirely in-dataset;
sole-observation, single-frame architectures in particular raise an
open question -- not addressed in that work or, to our knowledge,
anywhere else -- of whether relying on a single frame's local
appearance and context makes a model \emph{more} or \emph{less}
sensitive to the dataset-specific visual-context statistics we
localize as the dominant driver of cross-dataset sensitivity in
Section~\ref{sec:results-ablation}.

\subsection{Cross-Dataset and Robustness Evaluation}
\label{sec:related-robustness}

Gesnouin~et~al.~\cite{gesnouin2022} is, to our knowledge, the only prior
work to evaluate pedestrian crossing-intention predictors under a
cross-dataset protocol. Training eleven architectures spanning static
CNN, 3D-convolutional, pose-based, and RNN-based fusion designs
(including an SF-GRU-family model, SFRNN) on one of PIE, JAAD-behavior,
or JAAD-all and testing on the others, they find that every architecture
generalizes poorly -- performance is often statistically
indistinguishable from, or worse than, chance -- regardless of its
in-dataset ranking, and argue that classification metrics alone are
insufficient without accompanying uncertainty calibration (Expected and
Maximum Calibration Error). They further show that generic
ImageNet/Sports1M-pretrained backbones (VGG16, C3D, I3D) generalize
comparably to or better than crossing-intention-specific architectures
under domain shift, and that pretraining on data further from the target
domain improves calibration under shift. A citation-graph search over
all papers citing~\cite{gesnouin2022} through 2026 (spanning the
architectures in Section~\ref{sec:related-arch} onward) found none that
independently repeat a PIE$\leftrightarrow$JAAD cross-dataset
evaluation; a 2022 survey of the field cites it as the sole reference
for this open problem~\cite{fang2022survey}.

Two further papers from the same research group are directly relevant.
Achaji~et~al.~\cite{achaji2022bbox} show that bounding-box-only features
are unexpectedly sufficient for strong in-domain crossing-intention
accuracy on PIE, and hypothesize this may reflect a dataset-specific
bias rather than a genuinely informative signal; they test this via
synthetic-to-real transfer (a CARLA-simulated dataset to PIE) rather
than real-to-real cross-dataset transfer, and do not evaluate JAAD. This
is, to our knowledge, the closest prior evidence that box-trajectory
reliance can reflect a dataset-specific bias rather than a genuinely
transferable signal, a hypothesis we test directly with a real
cross-dataset, no-box ablation (Section~\ref{sec:results-geometry}); our
modality-zeroing results (Section~\ref{sec:results-ablation}), however,
localize cross-dataset sensitivity primarily to the visual
scene-context modality rather than to the box trajectory itself,
suggesting box-trajectory sufficiency and cross-dataset box-trajectory
fragility are related but not identical phenomena for SF-GRU.
Gilles~et~al.~\cite{gilles2022uncertainty} study cross-dataset
uncertainty for general autonomous-vehicle trajectory forecasting
(Argoverse, nuScenes, Interaction, and the Shifts benchmark) rather than
pedestrian crossing-intention classification on PIE/JAAD, but establish
that cross-dataset evaluation and uncertainty estimation are an active
concern in the broader AV-perception community, not one specific to
crossing-intention prediction.

Our work directly extends~\cite{gesnouin2022}'s finding to the current
generation of lightweight fusion architectures: we show the same
collapse persists for SF-GRU under corrected JAAD labels
(Section~\ref{sec:results-collapse}), that it is not resolved by
architectural refinements published since 2022
(Section~\ref{sec:related-arch}, none of which re-evaluate cross-dataset
performance), and that it is not explained by camera-geometry mismatch
or video-level split leakage, two plausible confounds neither prior
paper tests directly (Sections~\ref{sec:results-geometry}
and~\ref{sec:results-leakage}). Going beyond~\cite{gesnouin2022}'s
architecture-level comparison and~\cite{achaji2022bbox}'s in-domain
sufficiency finding, we localize the failure to a specific input
modality within a single architecture
(Section~\ref{sec:results-ablation}), apply a domain-adversarial and
contrastive mitigation (Section~\ref{sec:results-dann}), and use an
independent representation probe to show that mitigation's benefit is
dissociated from the representational invariance it is designed to
produce (Section~\ref{sec:results-probe}).

Uncertainty and calibration for pedestrian intention prediction has also
been studied outside the cross-dataset setting.
Zhang~et~al.~\cite{zhang2023trep} apply evidential deep learning within
a transformer to quantify prediction confidence under label noise.
Concurrent work on the PSI dataset~\cite{ashayer2026psi} incorporates a
variational bottleneck and a Mahalanobis-distance detector for
\emph{within}-dataset out-of-distribution sample detection and selective
prediction, improving accuracy by abstaining on the riskiest individual
test samples; this is a complementary but distinct problem from the
\emph{between}-dataset domain shift we study, and that work does not
evaluate PIE/JAAD or cross-dataset transfer.

\subsection{Positioning}
\label{sec:related-positioning}

Taken together, the literature on pedestrian crossing-intention
prediction has, since~\cite{gesnouin2022}, continued to report
in-dataset accuracy gains through richer fusion
(Section~\ref{sec:related-arch}) and lower inference cost
(Section~\ref{sec:related-efficiency}) without re-examining whether
those gains reflect genuine, transferable intention-recognition ability
or dataset-specific fit -- despite the open problem being visible enough
to be cited in a field survey~\cite{fang2022survey}. Our contribution is
to re-open that question empirically for the current architecture
generation, using a lightweight representative architecture (SF-GRU) for
which cross-dataset training and modality-level diagnosis is
computationally tractable, and to go beyond diagnosis to a partial,
mechanistically-understood fix: a domain-adversarial and contrastive
training recipe that reproducibly improves cross-dataset transfer in
both directions without achieving the literal representational
invariance it is nominally designed around
(Section~\ref{sec:results-probe}).
